%% ==============================================================================
%% content/03-methodologie.tex
%% Chapitre 3 : Styles, Macros Métier et Intégration Graphique (config/ & figures/)
%% ==============================================================================
\chapter{Styles, Macros Métier et Intégration Graphique}
\label{chap:styles-macros-graphisme}

\section{Organisation Modulaire du Dossier \code{config/}}
\label{sec:config-organisation}

Le répertoire \code{config/} sépare rigoureusement deux catégories d'instructions souvent confondues :
les paquets tiers d'une part, et les macros personnelles d'autre part.

\begin{itemize}
  \item \textbf{\code{config/packages.sty} :}
  déclare l'ensemble des dépendances logicielles externes nécessaires au document.
  Il configure leur comportement global (moteur de polices, langues, tableaux, hyperliens)
  dans le respect strict de l'ordre d'importation requis par \LaTeX{}.
  \item \textbf{\code{config/macros.sty} :}
  définit les outils métier spécifiques à la rédaction de votre ouvrage.
  On y trouve vos conventions mathématiques, vos raccourcis textuels
  et la personnalisation des environnements graphiques (\code{tcolorbox}).
\end{itemize}

Grâce à cette scission, si vous souhaitez réutiliser vos macros mathématiques dans un autre article,
il vous suffit d'importer \code{macros.sty} sans devoir importer toute l'infrastructure de la thèse.

\section{L'Art de Configurer les Paquets : \code{packages.sty}}
\label{sec:packages-detail}

Le fichier \code{packages.sty} met en \oe{}uvre plusieurs principes de composition moderne :

\subsection{Détection universelle du moteur typographique}
Le paquet \code{iftex} teste dynamiquement le compilateur utilisé :
\begin{itemize}
  \item Sous \textbf{Lua\LaTeX{}} ou \textbf{Xe\LaTeX{}}, il charge \code{fontspec} et active le crénage natif OpenType.
  \item Sous \textbf{pdf\LaTeX{}}, il charge de manière transparente le codage \code{[T1]\{fontenc\}}, l'encodage \code{[utf8]\{inputenc\}} et la fonte vectorielle \code{lmodern}.
\end{itemize}
Votre document compile ainsi sans friction sur n'importe quelle machine ou forge d'intégration continue.

\subsection{Microtypographie et lisibilité}
Le paquet \code{microtype} applique deux raffinements majeurs de l'imprimerie traditionnelle :
la saillie des marges (\textit{margin protrusion}, permettant aux tirets et à la ponctuation de déborder légèrement dans la marge blanche pour un alignement optique parfait)
et l'expansion subtile des caractères (\textit{font expansion}), qui élimine les césures excessives et les lignes lâches.

\subsection{Tableaux professionnels : la fin des filets verticaux}
Comme l'expliquent \textcite{bringhurst2004elements} et \textcite{mittelbach2023latex},
les tableaux typographiques de haute qualité n'utilisent **jamais** de filets verticaux.
Ces filets fragmentent le regard et gênent la lecture des données.
Notre configuration s'appuie sur le trio moderne :
\begin{enumerate}
  \item \textbf{Filets hiérarchisés :} le paquet \code{booktabs} (\code{\textbackslash toprule}, \code{\textbackslash midrule}, \code{\textbackslash bottomrule}),
  \item \textbf{Largeurs automatiques :} le paquet \code{tabularx} et sa colonne \code{X},
  \item \textbf{Cellules multilignes :} le paquet \code{makecell} pour les retours à la ligne internes.
\end{enumerate}

\section{Macros Sémantiques et Boîtes Graphiques : \code{macros.sty}}
\label{sec:macros-detail}

Le fichier \code{macros.sty} propose une boîte à outils prête à l'emploi.

\subsection{Raccourcis mathématiques rigoureux}
Pour éviter de ressaisir de longues commandes verbeuses, le fichier fournit :
\begin{itemize}
  \item Les vecteurs et matrices en gras droit : \code{\textbackslash vect\{x\}} ($\vect{x}$) et \code{\textbackslash mat\{M\}} ($\mat{M}$).
  \item Les ensembles de nombres usuels : \code{\textbackslash setR} ($\setR$), \code{\textbackslash setN} ($\setN$), \code{\textbackslash setC} ($\setC$).
  \item Les délimiteurs extensibles gérés par \code{mathtools} :
  \code{\textbackslash norm\{x\}} ($\norm{x}$),
  \code{\textbackslash abs\{x\}} ($\abs{x}$)
  et le produit scalaire \code{\textbackslash scal\{x\}\{y\}} ($\scal{x}{y}$).
  \item Le symbole différentiel droit conforme aux normes ISO : \code{\textbackslash diff} ($\diff x$).
\end{itemize}

\subsection{Environnements encadrés avec \code{tcolorbox}}
Trois types d'encadrés personnalisés sont pré-paramétrés :
\begin{enumerate}
  \item \code{\textbackslash begin\{callout\}[Titre]...} : encadré d'information générale ou de remarque méthodologique.
  \item \code{\textbackslash begin\{hypothesis\}[Titre]...} : boîte d'hypothèse scientifique, numérotée par chapitre.
  \item \code{\textbackslash begin\{definitionbox\}[Titre]...} : boîte de définition formelle avec en-tête plein.
\end{enumerate}

Voici une illustration en situation réelle de la boîte d'hypothèse :

\begin{hypothesis}[Découplage architectural strict]
\label{hyp:decouplage}
Tout composant graphique ou typographique peut être modifié ou remplacé
sans qu'aucune ligne des chapitres de contenu (\code{content/*.tex}) ne doive être éditée.
\end{hypothesis}

Grâce à l'intégration de \code{cleveref}, vous pouvez référencer automatiquement cet encadré dans le texte
via \code{\textbackslash cref\{hyp:decouplage\}}, ce qui produit automatiquement : \cref{hyp:decouplage}.

\section{Gestion des Figures Vectorielles : Standalone vs Externalisation TikZ}
\label{sec:figures-tikz}

\subsection{Le problème du recalcul systématique de TikZ}
Par défaut, à chaque compilation du manuscrit, le moteur \LaTeX{} recalcule et redessine l'ensemble
des figures TikZ et des tracés \code{pgfplots} à partir de zéro.
Sur un manuscrit de thèse comprenant 30 à 50 figures élaborées,
ce traitement peut ajouter plusieurs minutes à chaque cycle de compilation,
même si vous avez simplement corrigé une coquille typographique dans un paragraphe textuel.

\subsection{Pourquoi l'externalisation classique plante avec \code{out\_dir = 'build'} ?}
Pour accélérer ce processus, la bibliothèque native \code{\textbackslash usetikzlibrary\{external\}} propose
d'isoler chaque schéma dans un sous-processus d'arrière-plan via l'option \code{-shell-escape},
générant un PDF d'image réutilisé aux compilations ultérieures.
Toutefois, dès que le projet utilise un dossier de sortie isolé (\code{\$out\_dir = 'build'}),
deux conflits majeurs surgissent :
\begin{enumerate}
  \item \textbf{Conflit de répertoire :} TikZ cherche à écrire ses fichiers temporaires et ses sommes de contrôle
  dans un sous-dossier (ex. \code{tikz-cache/main-fig0.md5}). Si ce sous-dossier n'est pas préalablement créé
  à l'intérieur de \code{build/}, la compilation s'interrompt avec l'erreur :
  \code{! I can't write on file 'tikz-cache/main-fig0.md5'}.
  \item \textbf{Conflit de contexte de processus :} le processus principal travaille dans \code{build/},
  mais le sous-processus lancé par \code{\textbackslash write18} s'exécute à la racine.
  Le PDF généré par le sous-processus se retrouve alors égaré à la racine tandis que le moteur principal l'attend dans \code{build/}.
\end{enumerate}

\subsection{La solution reine en thèse : le paquet \code{standalone} (Option A)}
Pour s'affranchir définitivement de ces écueils et respecter les règles de sécurité
(évitant l'activation d'un \code{-shell-escape} non restreint),
notre architecture adopte le paquet \code{standalone}.

Chaque figure devient un fichier \LaTeX{} autonome dans le dossier \code{figures/tikz/} :
\begin{enumerate}
  \item \textbf{Édition ultra-rapide :} vous pouvez ouvrir \code{figures/tikz/schema-principe.tex}
  et le compiler isolément en moins d'une seconde dans votre éditeur pour ajuster vos n\oe{}uds.
  \item \textbf{Inclusion fluide :} dans le chapitre,
  l'instruction \code{\textbackslash includestandalone} extrait le tracé sans réexécuter de préambule :
\end{enumerate}

\begin{figure}[htbp]
  \centering
  \includestandalone[width=\textwidth]{figures/tikz/schema-principe}
  \caption[Schéma d'orchestration modulaire]{Schéma d'orchestration modulaire inclus via le paquet \code{standalone}}
  \label{fig:architecture-projet}
\end{figure}

Comme démontré sur le \cref{fig:architecture-projet},
le schéma hérite fidèlement de la palette chromatique de la thèse (\code{ThesisPrimary}, \code{ThesisSecondary})
grâce à la directive \code{\textbackslash providecolor} déclarée dans l'en-tête du fichier autonome.

\begin{callout}[Alternative : Précompilation en PDF]
Pour un temps de compilation de thèse strictement nul sur les figures stabilisées,
vous pouvez compiler le fichier \code{schema-principe.tex} une fois pour toutes
en \code{schema-principe.pdf}, puis l'appeler classiquement par \code{\textbackslash includegraphics}.
Vous obtenez ainsi une vélocité maximale sans compromettre la source vectorielle.
\end{callout}

\section{Sanctuarisation du Code Source Externe : le Dossier \code{listings/}}
\label{sec:listings-externe}

Une thèse ou un rapport technique intègre fréquemment des implémentations logicielles,
des scripts d'analyse ou des modules algorithmiques.
Le copier-coller de code directement dans les fichiers \code{.tex} constitue une mauvaise pratique notoire :
il détruit l'indentation, corrompt les caractères spéciaux
et empêche d'exécuter ou de tester le code en dehors du document.

Pour préserver une hygiène irréprochable, les codes sources sont placés
dans le répertoire dédié \code{listings/}.
Ils conservent leur extension native (\code{.py}, \code{.c}, \code{.rs}, \code{.jl}),
ce qui permet de les exécuter, de les profiler et de les tester unitairement avec vos outils habituels.
Dans le document \LaTeX{}, le code est inséré sans aucune duplication grâce au paquet \code{listings} :

\lstinputlisting[
  language=Python,
  caption={Pipeline d'optimisation modulaire importé directement depuis \code{listings/algorithm-pipeline.py}},
  label={lst:pipeline-python}
]{listings/algorithm-pipeline.py}

Comme illustré dans le \cref{lst:pipeline-python},
la mise en page hérite automatiquement de la typographie,
de la numérotation des lignes et de la palette chromatique (\code{ThesisPrimary}, \code{ThesisAccent}) définies dans \code{config/macros.sty}.

\section{Insertion Propre de Documents et Schémas PDF Externes}
\label{sec:insertion-pdf}

L'intégration de fichiers au format PDF dans une thèse répond à deux besoins fondamentalement distincts :
l'inclusion d'illustrations vectorielles d'une part,
et l'annexion de documents administratifs ou d'articles scientifiques complets d'autre part.

\subsection{Figures vectorielles : résolution infinie et rognage optique}
Pour les graphiques et schémas issus d'outils externes (Inkscape, Adobe Illustrator, Python ou MATLAB),
le format vectoriel PDF est la règle absolue :
contrairement aux images matricielles (\code{.png}, \code{.jpg}),
il ne pixellise jamais à l'agrandissement et garantit une impression d'une netteté parfaite.

Deux bonnes pratiques indispensables assurent une intégration professionnelle :
\begin{enumerate}
  \item \textbf{Omission de l'extension :} ne précisez jamais l'extension \code{.pdf} dans l'instruction :
  \begin{center}
    \code{\textbackslash includegraphics[width=0.8\textbackslash textwidth]\{figures/mon-schema\}}
  \end{center}
  \LaTeX{} choisira automatiquement le format vectoriel optimal.
  \item \textbf{Élimination des marges blanches parasites :} les logiciels de dessin exportent souvent une page A4 entière autour d'un petit schéma. Deux méthodes permettent d'isoler la boîte englobante exacte :
  \begin{itemize}
    \item \textbf{En amont (recommandé) :} exécuter l'outil en ligne de commande \code{pdfcrop} :
    \begin{center}
      \code{pdfcrop --margins 2 figures/mon-schema.pdf figures/mon-schema.pdf}
    \end{center}
    \item \textbf{Dans \LaTeX{} :} utiliser les clés de rognage \code{trim} et \code{clip} :
    \begin{center}
      \code{\textbackslash includegraphics[trim=10 20 10 20, clip, width=0.7\textbackslash textwidth]\{figures/mon-schema\}}
    \end{center}
  \end{itemize}
\end{enumerate}

\subsection{Documents multipages et publications en annexe : le paquet \code{pdfpages}}
Pour intégrer des pièces jointes officielles (conventions de recherche, attestations de jury, tirés-à-part d'articles publiés),
l'emploi de \code{\textbackslash includegraphics} est inadapté car il ne capture que la première page et tronque le document.

Le paquet \code{pdfpages} (\cref{sec:packages-detail}) répond à ce besoin via la commande \code{\textbackslash includepdf}.
Cependant, son usage direct présente un piège majeur :
par défaut, le style de page appliqué est \code{empty},
ce qui supprime les en-têtes et le numéro de page,
créant une rupture inacceptable dans la pagination continue.

Notre configuration dans \code{macros.sty} propose la commande dédiée :
\begin{center}
  \code{\textbackslash includepdfpaper[pages=\{1-3\}]\{figures/sample-annex-doc.pdf\}}
\end{center}

Cette commande applique deux réglages indispensables :
\begin{itemize}
  \item \code{pagecommand=\{\textbackslash thispagestyle\{scrheadings\}\}} :
  maintient le bandeau supérieur de chapitre et le folio (numéro de page) KOMA-Script au-dessus du document externe.
  \item \code{scale=0.82} :
  applique une réduction d'échelle subtile afin que le document annexé ne chevauche pas les en-têtes et respecte la grille typographique de votre ouvrage.
\end{itemize}
